\documentclass[10pt]{extarticle}
\usepackage[a4paper,margin=1.5cm]{geometry}
\usepackage[T1]{fontenc}
\usepackage{amsmath,amssymb}
\usepackage{graphicx}
\usepackage[export]{adjustbox}
\usepackage{multicol}
\usepackage{enumitem}
\usepackage{booktabs,array,tabularx}
\setlength{\columnsep}{0.7cm}
\setlength{\parindent}{0pt}
\setlength{\parskip}{3pt}
\pagestyle{empty}
\makeatletter
\newcommand{\qfigdef}[4]{\@namedef{qf@file@#1}{#2}\@namedef{qf@w@#1}{#3}\@namedef{qf@wide@#1}{#4}}
\newcommand{\qf@check}[1]{\@ifundefined{qf@file@#1}{\PackageError{score2}{Unknown figure id #1}{}}{}}
\newcommand{\qf@img}[4][]{\includegraphics[width=#2\dimexpr\@nameuse{qf@w@#4}pt\relax,max width=#3,max height=0.45\textheight,#1]{\@nameuse{qf@file@#4}}}
\newif\ifqf@full
\newif\ifqf@boxed
\newsavebox{\qf@box}
\newcommand{\qf@open}{\global\qf@fulltrue\par\global\qf@boxedtrue
  \setbox\qf@box\vbox\bgroup\hsize\textwidth\linewidth\textwidth}
\newcommand{\qf@emit}[1]{\par\egroup\global\qf@boxedfalse
  \ifdim\dimexpr\ht\qf@box+\dp\qf@box\relax<#1\textheight\noindent\usebox{\qf@box}\par\else\unvbox\qf@box\fi}
\newcommand{\qfullsplit}{\ifqf@boxed\qf@emit{0.85}\fi}
\newcommand{\qf@close}{\ifqf@boxed\qf@emit{0.45}\fi\global\qf@fullfalse}
\newcommand{\qfullbegin}{\end{multicols}\qf@open}
\newcommand{\qfullend}{\qf@close\begin{multicols}{2}}
\newcommand{\qfullstart}{\qf@open}
\newcommand{\qfulldone}{\qf@close}
\newcommand{\qfig}[2][1]{\qf@check{#2}%
  \ifqf@full\qf@blk{#1}{#2}\else
  \ifnum\@nameuse{qf@wide@#2}=1 \expandafter\@firstoftwo\else\expandafter\@secondoftwo\fi
  {\end{multicols}\begin{center}\qf@img{#1}{\textwidth}{#2}\end{center}\begin{multicols}{2}}%
  {\qf@blk{#1}{#2}}\fi}
\newcommand{\qf@blk}[2]{\par\smallskip{\centering\qf@img{#1}{\linewidth}{#2}\par}\smallskip}
\newcommand{\qfigi}[2][1]{\qf@check{#2}\qf@img[valign=c]{#1}{\linewidth}{#2}}
\makeatother
\newcommand{\papersec}[1]{\par\bigskip{\centering\Large\bfseries PAPER-#1\par}\medskip}
\newcommand{\partsec}[2]{\par\medskip\noindent{\large\bfseries PART-#1: #2}\par\smallskip}
\newcommand{\sect}[1]{\par\smallskip\noindent\textbf{#1}\par}
\newcommand{\qn}[1]{\par\smallskip\noindent\textbf{#1.}~}
\newcommand{\paraq}[2]{\par\smallskip\noindent\textbf{Paragraph for Questions #1 and #2}\par}
\newenvironment{qoptions}{\begin{itemize}[leftmargin=1.8em,itemsep=1pt,topsep=2pt,label={}]}{\end{itemize}}
\newcommand{\opt}[2]{\item[(#1)] #2}
% figure registry (generated by score2_tools.py finalize)
\qfigdef{p1_phy_s1_q02_1}{SRG (Enthusiast) - Test 29 (01 Oct 2023)-img/p1_phy_s1_q02_1.png}{264.5}{1}
\qfigdef{p1_phy_s1_q03_1}{SRG (Enthusiast) - Test 29 (01 Oct 2023)-img/p1_phy_s1_q03_1.png}{102.7}{0}
\qfigdef{p1_phy_s1_q08_1}{SRG (Enthusiast) - Test 29 (01 Oct 2023)-img/p1_phy_s1_q08_1.png}{78.2}{0}
\qfigdef{p1_phy_s2_q10_1}{SRG (Enthusiast) - Test 29 (01 Oct 2023)-img/p1_phy_s2_q10_1.png}{76.2}{0}
\qfigdef{p1_phy_s2_q10_2}{SRG (Enthusiast) - Test 29 (01 Oct 2023)-img/p1_phy_s2_q10_2.png}{92.5}{0}
\qfigdef{p1_phy_s2_q10_3}{SRG (Enthusiast) - Test 29 (01 Oct 2023)-img/p1_phy_s2_q10_3.png}{94.2}{0}
\qfigdef{p1_phy_s2_q10_4}{SRG (Enthusiast) - Test 29 (01 Oct 2023)-img/p1_phy_s2_q10_4.png}{93.3}{0}
\qfigdef{p1_phy_s3_q02_1}{SRG (Enthusiast) - Test 29 (01 Oct 2023)-img/p1_phy_s3_q02_1.png}{83.3}{0}
\qfigdef{p1_phy_s3_q07_1}{SRG (Enthusiast) - Test 29 (01 Oct 2023)-img/p1_phy_s3_q07_1.png}{127.2}{0}
\qfigdef{p1_chem_s1_q06_1}{SRG (Enthusiast) - Test 29 (01 Oct 2023)-img/p1_chem_s1_q06_1.png}{291.1}{1}
\qfigdef{p1_chem_s2_q09_1}{SRG (Enthusiast) - Test 29 (01 Oct 2023)-img/p1_chem_s2_q09_1.png}{94.1}{0}
\qfigdef{p1_chem_s2_q09_2}{SRG (Enthusiast) - Test 29 (01 Oct 2023)-img/p1_chem_s2_q09_2.png}{86.1}{0}

\begin{document}
\begin{center}{\Large\bfseries SRG (Enthusiast) --- Test 29 (01 Oct 2023)}\end{center}
\begin{multicols}{2}
\partsec{1}{PHYSICS}
\sect{Section-I(i)}
\qn{1} The reactance of a circuit is zero. It is possible that the circuit contains:
\begin{qoptions}
\opt{A}{An inductor and a capacitor}
\opt{B}{An inductor but no capacitor}
\opt{C}{A capacitor but no inductor}
\opt{D}{Neither an inductor nor a capacitor}
\end{qoptions}
\qfullbegin
\qn{2} Current growth in two L-R circuits(b) and (c) as shown in fig.(a). Let $L_1$, $L_2$, $R_1$ and $R_2$ be the corresponding values in two circuits, Then:
\qfig{p1_phy_s1_q02_1}
\qfullsplit
\begin{qoptions}
\opt{A}{$R_1>R_2$}
\opt{B}{$R_1=R_2$}
\opt{C}{$L_1>L_2$}
\opt{D}{$L_1<L_2$}
\end{qoptions}
\qfullend
\qn{3} An inductor $L$, a resistance $R$ and two identical bulbs $B_1$ and $B_2$ are connected to a battery through a switch $S$ as shown in the figure. The resistance of coil having inductance $L$ is also $R$. Which of the following statement gives the correct description of the happenings when the switch $S$ is closed?
\qfig{p1_phy_s1_q03_1}
\begin{qoptions}
\opt{A}{The bulb $B_2$ lights up earlier than $B_1$ and finally both the bulbs shine equally bright.}
\opt{B}{$B_1$ light up earlier and finally both the bulbs acquire equal brightness.}
\opt{C}{$B_2$ lights up earlier and finally $B_1$ shines brighter than $B_2$.}
\opt{D}{$B_1$ and $B_2$ light up together with equal brightness all the time.}
\end{qoptions}
\qn{4} A series R-C circuit is connected to AC voltage source. Consider two cases; (A) when $C$ is without a dielectric medium and (B) when $C$ is filled with dielectric of constant 4. The current $I_R$ through the resistor and voltage $V_C$ across the capacitor are compared in the two cases. Which of the following is/are true?
\begin{qoptions}
\opt{A}{$I_R^A>I_R^B$}
\opt{B}{$I_R^A<I_R^B$}
\opt{C}{$V_C^A>V_C^B$}
\opt{D}{$V_C^A<V_C^B$}
\end{qoptions}
\qn{5} A horizontal plank has a rectangular block placed on it. The plank starts oscillating vertically and simple harmonically with an amplitude of 40 cm. The block just loses contact with the plank when the letter is at momentary rest. Then
\begin{qoptions}
\opt{A}{The period of oscillation is $\left(\dfrac{2\pi}{5}\right)$}
\opt{B}{The block weighs double its weight. when the plank is at one of the positions of momentary rest.}
\opt{C}{The block weighs 0.5 times its weight on the plank halfway up from mean position}
\opt{D}{The block weighs 1.5 times its weight on the plank halfway down from mean position}
\end{qoptions}
\qn{6} For a body executing SHM with amplitude $A$, time period $T$, maximum velocity $V_{\max}$ and phase constant zero. Which of the following statements are correct for $0\le t\le\dfrac T4$ ($y$ is displacement from mean position)?
\begin{qoptions}
\opt{A}{At $y=(A/2)$, $v>(v_{\max}/2)$}
\opt{B}{For $v=(v_{\max}/2)$, $y>(A/2)$}
\opt{C}{For $t=(T/8)$, $y>(A/2)$}
\opt{D}{For $y=(A/2)$, $t<(T/8)$}
\end{qoptions}
\qn{7} The potential energy of a particle of mass 0.1 kg, moving along the $x$-axis, is given by $U=5x\,(x-4)$ J, where $x$ is in meters. It can be concluded that
\begin{qoptions}
\opt{A}{The particle is acted by a constant force}
\opt{B}{The speed of the particle is maximum at $x=2$m}
\opt{C}{The particle executes SHM}
\opt{D}{The period of oscillation of the particles is $(\pi/5)$ sec}
\end{qoptions}
\qn{8} A point $P$ is located at a distance of 3 m from a point light source of power 100 W. Assume the efficiency of the bulb is 2.5\%
\qfig{p1_phy_s1_q08_1}
\begin{qoptions}
\opt{A}{The intensity at point $P$ is 0.022 W/m$^2$}
\opt{B}{The intensity at point $P$ is 0.011 W/m$^2$}
\opt{C}{The intensity of electric field at point $P$ is 0.022 W/m$^2$}
\opt{D}{The intensity of magnetic field at point $P$ is 0.011 W/m$^2$}
\end{qoptions}
\sect{Section-I(ii)}
\qn{9} Match the column I with column II\\
In simple harmonic motion: $x=1.0\sin[12\pi t]$ and mass of particle executing SHM, $m=1/4$ kg
\begin{center}\small\renewcommand{\arraystretch}{1.5}
\begin{tabular}{|l|p{0.5\linewidth}||l|l|}\hline
\multicolumn{2}{|c||}{\textbf{Column A}} & \multicolumn{2}{c|}{\textbf{Column B}}\\\hline
(A) & Frequency with which kinetic energy oscillates=..... & (p) & $\dfrac1{12}$\\\hline
(B) & Speed of particle is maximum at time=... & (q) & $18\pi^2$\\\hline
(C) & Maximum potential energy=... & (r) & 12\\\hline
(D) & Force constant $k$=... & (s) & $36\pi^2$\\\hline
\end{tabular}
\end{center}
\begin{qoptions}
\opt{A}{$A\to p;\ B\to r;\ C\to s;\ D\to q$}
\opt{B}{$A\to r;\ B\to p;\ C\to q;\ D\to s$}
\opt{C}{$A\to s;\ B\to p;\ C\to r;\ D\to p$}
\opt{D}{$A\to q;\ B\to p;\ C\to s;\ D\to r$}
\end{qoptions}
\qn{10} In column I, variation of current $I$ with time $t$ is given in figures. In column II root means square current $I$ rms, and average current is given. Match the column-I with corresponding quantities given in column-II
\begin{center}\small\renewcommand{\arraystretch}{1.3}
\begin{tabular}{|l|c||l|p{0.30\linewidth}|}\hline
\multicolumn{2}{|c||}{Column-II} & \multicolumn{2}{c|}{Column-II}\\\hline
(P) & \qfigi{p1_phy_s2_q10_1} & (1) & $i_{\rm rms}=\dfrac{i_0}{\sqrt3}$\\\hline
(Q) & \qfigi{p1_phy_s2_q10_2} & (2) & Average current for positive half cycle is $i_0$\\\hline
(R) & \qfigi{p1_phy_s2_q10_3} & (3) & Average current for positive half cycle is $\dfrac{i_0}{2}$\\\hline
(S) & \qfigi{p1_phy_s2_q10_4} & (4) & Full cycle average current is zero\\\hline
\end{tabular}
\end{center}
\begin{qoptions}
\opt{A}{$P\to1,4;\ Q\to1,4;\ R\to3,4;\ S\to2$}
\opt{B}{$P\to4;\ Q\to2,3,4;\ R\to2,4;\ S\to2$}
\opt{C}{$P\to4;\ Q\to1,3,4;\ R\to1,4;\ S\to2,4$}
\opt{D}{$P\to4;\ Q\to1,3,4;\ R\to2,4;\ S\to2$}
\end{qoptions}
\sect{Section-II}
\qn{1} A closed circuit consists of a source of constant emf $E$ and a choke coil of inductance $L$ connected in series. The active resistance of the whole circuit is equal to $R$. It is in steady state. At the moment $t=0$ the choke coil inductance was decreased abruptly 4 times. The current in the circuit as a function of time $t$ is $\dfrac ER\left[1+xe^{-4tR/L}\right]$. Find out value of $x$.
\qn{2} In the circuit diagram shown in the figure the switches $S_1$ and $S_2$ are closed to time $t=0$. After time $t=(0.1)\,\ell$n2, switch $S_2$ is opened. The current in the circuit at time, $t=(0.2)\,\ell$n2 sec is equal to $\dfrac x{32}$ amp. Find out value of $x$.
\qfig{p1_phy_s3_q02_1}
\qn{3} A solenoid of length 1 m, area of cross-section 4.0 cm$^2$ and having 4000 turns is placed inside another solenoid of 2000 turns having a cross-sectional area 6 cm$^2$ and length 2 m. The mutual inductance between the solenoids is $x\pi\times10^{-6}$H. Find out value of $x$.
\qn{4} A series LCR circuit containing a resistance of 120 ohm has angular resonance frequency $4\times10^5$ rad s$^{-1}$. At resonance, the voltage across resistance and inductance are 60 V and 40 V respectively. At frequency the current in the circuit lags the voltage by $45^\circ$ is equal to $x\times10^5$ rad/sec. Find value of $x$
\qn{5} An LCR circuit has $L=10$ mH, $R=150\,\Omega$ and $C=1\,\mu$F connected in series to a source of $150\sqrt2\cos\omega t$ volt. At a frequency that is 50\% of the resonant frequency, calculate the average power(in watt) dissipated per cycle
\qn{6} A particle of mass 1.5 kg moves along the $x$-axis in force field, Its potential energy is given by $V(x)=2x^3-9x^2+12x$, where all quantities are written in S.I units. There is an equilibrium point on the $x$-axis at which if displaced slightly it performs SHM. If the angular frequency of oscillation is $x$ rad/s. Find $x$.
\qn{7} In the figure shown, mass of the plank is $m$ and that of the solid cylinder is $8m$. Springs are light. The planks is slightly displaced from equilibrium and then released. Find the period of small oscillations(in seconds) of the plank. There is no slipping at any contact point. The ratio of the mass of the plank and stiffness of the spring i.e $\dfrac mK=\dfrac2{\pi^2}$.
\qfig{p1_phy_s3_q07_1}
\qn{8} A 15.0 mW helium-neon laser($\lambda=632.8$ nm) emits a beam of circular cross section with a diameter of 2 mm. If the total energy is contained in a 1 m length of the beam is $x\times10^{-12}$J then find $x$.
\partsec{2}{CHEMISTRY}
\sect{Section-I(i)}
\qn{1} Select the correct statements regarding salt solution of weak acid \& strong base
\begin{qoptions}
\opt{A}{Solution will be neutral}
\opt{B}{Degree of hydrolysis increase with increase in dilution}
\opt{C}{Shows anionic hydrolysis}
\opt{D}{Phenolphthalein is suitable indicator during titration between weak acid \& strong base}
\end{qoptions}
\qn{2} For the cell $\mathrm{Cu_{(s)}|CuSO_{4(aq)}\,\|\,HCl_{(aq)}\,|H_{2(g)}|\,Pt_{(s)}}$ $E_{\text{cell}}$ can be increased by -
\begin{qoptions}
\opt{A}{By adding $\mathrm{NH_3}$ in anodic half cell}
\opt{B}{Increasing partial pressure of $\mathrm{H_2}$ gas}
\opt{C}{By diluting $\mathrm{CuSO_4}$ solution}
\opt{D}{By diluting HCl solution}
\end{qoptions}
\qn{3} Based on above data select correct statement under standard condition -\\
$E^\circ_{A^+/A}=0.5$V ; $E^\circ_{B^+/B}=0.2$V\\
$E^\circ_{C^+/C}=-0.3$V ; $E^\circ_{D^+/D}=-0.1$V
\begin{qoptions}
\opt{A}{Best oxidising agent is $\mathrm{A^+}$}
\opt{B}{Best reducing agent is D}
\opt{C}{C can displace D from it's solution}
\opt{D}{In galvanic cell obtained by electrodes A \& B current flow from B to A in external circuit}
\end{qoptions}
\qn{4} Buffer solution A of a weak monoprotic acid and its sodium salt in the concentration ratio $x:y$ has $\mathrm{p^H}=(\mathrm{pH})_1$. Buffer solution B of the same acid and its sodium salt in the concentration ratio $y:x$ has $\mathrm{p^H}=(\mathrm{pH})_2$. If $(\mathrm{pH})_2-(\mathrm{pH})_1=1$ unit and $(\mathrm{pH})_1+(\mathrm{pH})_2=9.5$ units, then
\begin{qoptions}
\opt{A}{$\mathrm{pK_a}=4.75$}
\opt{B}{$\dfrac xy=2.36$}
\opt{C}{$\dfrac xy=3.162$}
\opt{D}{$\mathrm{pK_a}=5.25$}
\end{qoptions}
\qn{5} At equilibrium state, which of the following features are correct?\\
$\underset{\text{(solid)}}{\mathrm{NH_2COONH_4}}\rightleftharpoons\underset{\text{(gas)}}{2\mathrm{NH_3}}+\underset{\text{(gas)}}{\mathrm{CO_2}}$
\begin{qoptions}
\opt{A}{Total pressure is 1.5 times that of partial pressure of $\mathrm{NH_3}$.}
\opt{B}{$K_P$ value is 4 times that of $(P_{\mathrm{CO_2}})^3$.}
\opt{C}{$P_{\mathrm{NH_3}}=\dfrac12P_{\mathrm{CO_2}}$}
\opt{D}{$\mathrm{CO_2}$ contributes 33.3\% of equilibrium pressure}
\end{qoptions}
\qfullbegin
\qn{6} Consider the following diagrams, for the three different reactions 1, 2 and 3 where, for each reaction:\\
$[\mathrm A]\to[\mathrm I]$ First order rate constant, $k_1$\\
$[\mathrm I]\to[\mathrm P]$ First order rate constant, $k_2$\\
$\mathrm A\xrightarrow{k_1}\mathrm I\xrightarrow{k_2}\mathrm P$
\qfig{p1_chem_s1_q06_1}
\qfullsplit
\begin{qoptions}
\opt{A}{The ratio of $k_1/k_2$ will be highest in the reaction (2)}
\opt{B}{$t_{\max}$ is directly proportional to $\log(k_2/k_1)$}
\opt{C}{The ratio of $\log(k_2/k_1)$ will follow the order: (2) $<$ (1) $<$ (3)}
\opt{D}{For a sequential reaction, in which $k_1=2k_2=0.1\mathrm{s^{-1}}$, the time taken for the maximum concentration of [I] will be approx. 14 sec}
\end{qoptions}
\qfullend
\qn{7} For a reaction mechanism\\
$\mathrm{A_2}\underset{k_2}{\overset{k_1}{\rightleftharpoons}}2\mathrm A$ (fast)\\
: $\mathrm A+\mathrm B\xrightarrow{k_3}$ Product (slow)\\
Which of the following relations are correct\,?
\begin{qoptions}
\opt{A}{Rate $=k_3\left(\dfrac{k_1}{k_2}\right)^{1/2}[\mathrm{A_2}]^{1/2}[\mathrm B]$}
\opt{B}{Rate $=k_3\sqrt{\dfrac{k_1}{k_2}\times k_{\text{eq}}}\,[\mathrm{A_2}]^{1/2}[\mathrm B]$}
\opt{C}{$-\dfrac{d[\mathrm B]}{dt}=k_3k_{\text{eq}}^{1/2}[\mathrm{A_2}]^{1/2}[\mathrm B]$}
\opt{D}{$-\dfrac{d[\mathrm A]}{dt}=k_2\left(\dfrac{k_1}{k_2}\right)^{1/2}[\mathrm{A_2}]^{1/2}[\mathrm B]$}
\end{qoptions}
\qn{8} Ozone decomposes according to the equation $2\mathrm{O_3(g)}\to3\mathrm{O_2(g)}$. Mechanism of the reaction is\\
Step I : $\mathrm{O_3(g)}\rightleftharpoons\mathrm{O_2(g)}+\mathrm{O(g)}$ (fast)\\
Step II : $\mathrm{O_3(g)}+\mathrm{O_2(g)}\to2\mathrm{O_2(g)}$ (slow)\\
Which of the following is correct?
\begin{qoptions}
\opt{A}{For step I, molecularity is 2}
\opt{B}{For step II, molecularity is 1}
\opt{C}{Step II is rate determining step}
\opt{D}{Rate law expression for the overall reaction is $-\dfrac d{dt}[\mathrm{O_3}]=k[\mathrm{O_3}]^2[\mathrm{O_2}]^{-1}$}
\end{qoptions}
\sect{Section-I(ii)}
\qn{9} Match the following
\begin{center}\footnotesize\setlength\tabcolsep{3pt}
\begin{tabular}{|l|p{0.30\linewidth}|l|p{0.30\linewidth}|}\hline
\multicolumn{2}{|c|}{\textbf{Column-I}} & \multicolumn{2}{c|}{\begin{tabular}{c}\textbf{Column-II}\\\textbf{(Suitable indicator at end point)}\end{tabular}}\\\hline
A & Strong acid vs strong base titration & P & Methyl orange (3.1-4.4)\\\hline
B & Weak acid vs strong titration & Q & Methyl red (4.2-6.3)\\\hline
C & Strong acid vs weak base titration & R & Phenolphthalein (8.3-10)\\\hline
D & Weak acid vs weak base titration & S & None of the above mentioned indicator is suitable\\\hline
 & & T & All of the above mentioned indicators are suitable\\\hline
\end{tabular}
\end{center}
\begin{qoptions}
\opt{A}{$A\to$ PQRT; $B\to$ R; $C\to$ PQ; $D\to$ S}
\opt{B}{$A\to$ PQR; $B\to$ S; $C\to$ PQ; $D\to$ R}
\opt{C}{$A\to$ PRT; $B\to$ PQ; $C\to$ S; $D\to$ R}
\opt{D}{$A\to$ PSRT; $B\to$ R; $C\to$ PS; $D\to$ Q}
\end{qoptions}
\qn{10} Match Matrix
\begin{center}\footnotesize\setlength\tabcolsep{3pt}
\begin{tabular}{|l|p{0.30\linewidth}|l|p{0.30\linewidth}|}\hline
\multicolumn{2}{|c|}{\textbf{Column - I}} & \multicolumn{2}{c|}{\textbf{Column - II}}\\\hline
A & Ester $+$ NaOH $\to$ Alcohol $+$ Salt & P & $k=Ae^{-E_a/RT}$\\\hline
B & $\mathrm{2H_2O_2\to2H_2O}$ $+$ $\mathrm{O_2}$ & Q & \qfigi{p1_chem_s2_q09_1}\\\hline
C & $\mathrm{2NH_3}\xrightarrow{\mathrm{Au}}\mathrm{N_2+3H_2}$ & R & \qfigi{p1_chem_s2_q09_2}\\\hline
D & Sucrose $+$ $\mathrm{H_2O}\to$ Glucose $+$ Fructose & S & $t_{1/2}=x$ min at a constant temperature any time of the reaction.\\\hline
\end{tabular}
\end{center}
\begin{qoptions}
\opt{A}{$A\to$ PQ; $B\to$ PRS; $C\to$ QRS; $D\to$ P}
\opt{B}{$A\to$ PR; $B\to$ PRS; $C\to$ Q; $D\to$ PRS}
\opt{C}{$A\to$ RS; $B\to$ Q; $C\to$ PS; $D\to$ RQ}
\opt{D}{$A\to$ PQ; $B\to$ QRS; $C\to$ SR; $D\to$ P}
\end{qoptions}
\sect{Section-II}
\qn{1} How many litre of 0.3 M NaOH solution should be added to a solution containing 0.5 moles of HCl in hydrogen electrode, to make its Oxidation potential $=0.042$ V at $25^\circ$C. $\left(\log5=0.7,\ \dfrac{2.303RT}{F}=0.06\text{V}\right)$
\qn{2} The solubility of $\mathrm{Pb(OH)_2}$ in water is $5\times10^{-6}$M. Calculate the pH of buffer solution in which solubility of $\mathrm{Pb(OH)_2}$ is $5\times10^{-4}$ is mol litre$^{-1}$.
\qn{3} A certain mixture of $\mathrm{H_2SO_4}$ and weak acid HA is 0.2 M in each of the acids and 20 mL of this solution is titrated against 20 mL of 0.5 M NaOH, the resultant pH of solution can be $K_a$(HA) $=10^{-5}$
\qn{4} When 3.06 g of solid $\mathrm{NH_4HS}$ is introduced into a two-liter evacuated flask at $27^\circ$C, 30\% of the solid decomposes into gaseous ammonia and hydrogen sulphide
Calculate $K_P$ for the reaction at $27^\circ$C $\left(R=0.082\ \text{lit atm mol}^{-1}\text{k}^{-1}\right)$
\qn{5} Consider a first order gaseous reaction at constant temperature and constant volume\\
$2\mathrm A\to\mathrm B(g)+\mathrm C(s)$\\
Initial pressure of $\mathrm A(g)$ is 75 mm Hg. After 10 minute the pressure of B gas was found to be 18.75 mm Hg. Calculate the time at which the pressure of A gas becomes 9.375 mm Hg? (in min).
\qn{6} $2\mathrm A(g)+\mathrm B(g)\to\mathrm C(g)$, these data were obtained when this reaction was studied
\begin{center}\small
\begin{tabular}{|c|c|c|}\hline
$[\mathrm A]$, M & $[\mathrm B]$, M & Rate of reaction $\left(\mathrm{M\,s^{-1}}\right)$\\\hline
0.05 & 0.01 & $2\times10^{-5}$\\\hline
0.25 & 0.01 & $10^{-4}$\\\hline
0.05 & 0.02 & $8\times10^{-4}$\\\hline
\end{tabular}
\end{center}
Calculate rate constant for this reaction (in the given units)
\qn{7} The standard gibbs free energy change $\left(\text{in kJ mol}^{-1}\right)$ for the following reaction is \_\_\_\_ (up to two decimal places)\\
$\mathrm{Sn(s)+Pb^{2+}=Sn^{2+}+Pb(s)}$\\
Given,\\
$E^0\left(\mathrm{Pb^{2+}/Pb}\right)-0.126$V,\\
$E^0\left(\mathrm{Sn^{2+}/Sn}\right)=-0.14$V\\
take $F=96500$C
\qn{8} The time required for 10\% completion of first order reaction at 298 K is equal to that required for its 25\% completion at 308 K. What is value of $\dfrac{k_{308}}{k_{208}}$\,?\\
{[}Given : $\log2=0.3$, $\log3=0.4771$, $\ell$n(10) $=2.303$]
\partsec{3}{MATHEMATICS}
\sect{Section-I(i)}
\qn{1} For $a,b>0$, the value of \par\smallskip\noindent\resizebox{\linewidth}{!}{$\displaystyle\int_0^1\dfrac{\log\left(\dfrac{(x+a)^{(x+a)}}{(x+b)^{(x+b)}}\right)}{(x+a)(x+b)\log(x+a)\log(x+b)}\,dx$}\par\smallskip\noindent is equal to
\begin{qoptions}
\opt{A}{$\log\left(\dfrac{1+a}{1+b}\right)\dfrac{\log a}{\log b}$}
\opt{B}{$\dfrac{\log(1+a)\log a}{\log(1+b)\log b}$}
\opt{C}{$\log\left(\dfrac{\log(1+a)\log a}{\log(1+b)\log b}\right)$}
\opt{D}{$\log\left(\dfrac{\log(1+b)\log a}{\log(1+a)\log b}\right)$}
\end{qoptions}
\qn{2} For $k\in N$, let $b_k=\displaystyle\int_1^e\dfrac1{(k+\log x)(k+1+\log x)}\dfrac{dx}x$ and $a_n=\displaystyle\sum_{k=1}^nb_k$. Then
\begin{qoptions}
\opt{A}{$b_k=\log\left(\dfrac{(k+1)^2}{k(k+2)}\right)$}
\opt{B}{$a_n=\log\left(\dfrac{2n+2}{n+2}\right)$}
\opt{C}{$\lim\limits_{n\to\infty}a_n=\log2$}
\opt{D}{$\lim\limits_{n\to\infty}a_n<1$}
\end{qoptions}
\qn{3} For $a>0$, let $I(a)=\displaystyle\int_0^a\dfrac{dx}{1+e^{a^2-2ax}}$ and let $0<\theta<\dfrac\pi2$, then
\begin{qoptions}
\opt{A}{$I'(a)=\ln\sqrt e$}
\opt{B}{$I(\sin\theta)+I(\cos\theta)\le\sin\left(\dfrac\pi4\right)$}
\opt{C}{$I\left(\sin^2\theta\right)+I\left(\cos^2\theta\right)=\dfrac12$}
\opt{D}{$I\left(2a-a^2\right)\le\dfrac12$}
\end{qoptions}
\qn{4} Let $f$ be a nonnegative function defined on the interval $\left[0,\dfrac\pi2\right]$. If\\
$\displaystyle\int_0^x\left(f'(t)-\sin2t\right)dt=\int_x^0f(t)\tan t\,dt$ and $f(0)=1$, then
\begin{qoptions}
\opt{A}{$f(0)=1$ is the absolute maximum value of $f$}
\opt{B}{$f(0)=1$ is the minimum value of $f$}
\opt{C}{$\displaystyle\int_0^{\pi/2}y\,dx=3-\dfrac\pi2$}
\opt{D}{$f\left(\dfrac\pi3\right)=1$}
\end{qoptions}
\qn{5} Let $f:[0,1]\to R$ be continuous and differentiable function such that for any $x,y\in[0,1]$ $xf(y)+yf(x)\le1$, then which of the following options is/are correct?
\begin{qoptions}
\opt{A}{$\displaystyle\int_0^1f(x)\,dx\le\dfrac\pi4$}
\opt{B}{$\displaystyle\int_0^1x^2f'(x)\,dx\ge f(1)-1$}
\opt{C}{$\displaystyle\int_0^1xf'(x)\,dx\ge f(1)-\dfrac\pi4$}
\opt{D}{$\displaystyle\int_0^1f(x)\,dx\le\dfrac\pi8$}
\end{qoptions}
\qn{6} If $f(x)=\lim\limits_{m\to\infty}\dfrac1m\displaystyle\sum_{k=0}^m\left|x-\dfrac km\right|$; then
\begin{qoptions}
\opt{A}{$f(0)=\dfrac12$}
\opt{B}{$f(1)=\dfrac12$}
\opt{C}{$f(-1)=\dfrac32$}
\opt{D}{$f(2)=\dfrac32$}
\end{qoptions}
\qn{7} Let $f(x)$ is a differentiable function such that $f'(x)=f(x)+\displaystyle\int_0^2f(x)\,dx$ if $f(0)=\dfrac{4-e^2}3$, then
\begin{qoptions}
\opt{A}{$f(x)=e^x-\left(\dfrac{e^2-1}3\right)$}
\opt{B}{$f(x)=e^x+\left(\dfrac{e^2-1}3\right)$}
\opt{C}{$f(0)<0$}
\opt{D}{$f(x)$ has atleast one root between 0 \& 1}
\end{qoptions}
\qn{8} A differentiable function satisfies $f(x)=\displaystyle\int_0^x\{f(t)\cos t-\cos(t-x)\}\,dt$. Which is of the following hold good?
\begin{qoptions}
\opt{A}{$f(x)$ has a minimum value $1-e$}
\opt{B}{$f(x)$ has a maximum value $1-e^{-1}$}
\opt{C}{$f''\left(\dfrac\pi2\right)=e$}
\opt{D}{$f'(0)=1$}
\end{qoptions}
\sect{Section-I(ii)}
\qn{9} \begin{center}\footnotesize\setlength\tabcolsep{2pt}\renewcommand{\arraystretch}{1.4}
\begin{tabular}{|c|p{0.55\linewidth}|c|c|}\hline
\multicolumn{2}{|c|}{\textbf{Column I}} & \multicolumn{2}{c|}{\textbf{Column II}}\\\hline
(A) & The function $f(x)=\dfrac{e^{x\cos x}-1-x}{\sin x^2}$ is not defined at $x=0$. The value of $f(0)$ so that $f$ is continuous at $x=0$ is & (P) & $-1$\\\hline
(B) & The value of the definite integral $\displaystyle\int_0^1\dfrac{dx}{\sqrt x+\sqrt[3]x}$ equals $a+b\ln2$, where $a$ and $b$ are integers then $(a+b)$ equals & (Q) & 0\\\hline
(C) & Given $e^n\displaystyle\int_0^n\dfrac{\sec^2\theta-\tan\theta}{e^\theta}d\theta=1$ then the value of $\tan(n)$ is equal to & (R) & $\dfrac12$\\\hline
(D) & Let $a_n=\displaystyle\int_{\frac1{n+1}}^{\frac1n}\tan^{-1}(nx)\,dx$ and $b_n=\displaystyle\int_{\frac1{n+1}}^{\frac1n}\sin^{-1}(nx)\,dx$ then $\lim\limits_{n\to\infty}\dfrac{a_n}{b_n}$ has the value equal to & (S) & 1\\\hline
\end{tabular}
\end{center}
\begin{qoptions}
\opt{A}{$A\to R;\ B\to P;\ C\to S;\ D\to R$}
\opt{B}{$A\to P;\ B\to Q;\ C\to S;\ D\to R$}
\opt{C}{$A\to R;\ B\to P;\ C\to Q;\ D\to S$}
\opt{D}{$A\to S;\ B\to Q;\ C\to R;\ D\to P$}
\end{qoptions}
\qn{10} Match the columns.\par\smallskip
\noindent\textbf{Column I}\par
\noindent(A)\ The value of the definite integral, $\displaystyle\int_{\theta_1}^{\theta_2}\dfrac{d\theta}{1+\tan\theta}=\dfrac{501\pi}K$ where $\theta_2=\dfrac{1003\pi}{2008}$ and $\theta_1=\dfrac\pi{2008}$. The value of $K$ equals\par\smallskip
\noindent(B)\ Suppose that the function $f,g,f'$ and $g'$ are continuous over $[0,1]$, $g(x)\ne0$ for $x\in[0,1]$, $f(0)=0$, $g(0)=\pi$, $f(1)=\dfrac{2009}2$ and $g(1)=1$. The value of the definite integral, \par\smallskip\noindent\resizebox{\linewidth}{!}{$\displaystyle\int_0^1\dfrac{f(x)\cdot g'(x)\left\{g^2(x)-1\right\}+f'(x)\cdot g(x)\left\{g^2(x)+1\right\}}{g^2(x)}dx$}\par\smallskip\noindent is equal to\par\smallskip
\noindent(C)\ The value of the integral \par\smallskip\noindent\resizebox{\linewidth}{!}{$\displaystyle\int_0^1\left(\sqrt[2007]{1-x^{2008}}-\sqrt[2008]{1-x^{2007}}\right)dx$}\par\smallskip\noindent, is equal to\par\smallskip
\noindent(D)\ $I=\displaystyle\int_2^3\left((x-1)^3+(4-x)^3+x\right)\cos\pi x\,dx$ then $\left|50\pi^2I\right|=$\par\smallskip
\noindent\textbf{Column II}\par
\begin{center}\small\begin{tabular}{cccc}(P) 0 & (Q) 100 & (R) 2008 & (S) 2009\end{tabular}\end{center}
\begin{qoptions}
\opt{A}{$A\to P;\ B\to S;\ C\to R;\ D\to Q$}
\opt{B}{$A\to R;\ B\to S;\ C\to P;\ D\to Q$}
\opt{C}{$A\to S;\ B\to R;\ C\to P;\ D\to Q$}
\opt{D}{$A\to Q;\ B\to S;\ C\to R;\ D\to P$}
\end{qoptions}
\sect{Section-II}
\qn{1} Let $m$ be the minimum possible value of $\log_3\left(3^{y_1}+3^{y_2}+3^{y_3}\right)$, where $y_1,y_2,y_3$ are real numbers for which $y_1+y_2+y_3=9$. Let $M$ be the maximum possible value of $(\log_3x_1+\log_3x_2+\log_3x_3)$, where $x_1,x_2,x_3$ are positive real numbers for which $x_1+x_2+x_3=9$. Then the value of $\log_2\left(m^3\right)+\log_3\left(M^2\right)$ is \_\_\_\_\_\_\_\_.
\qn{2} Find the number of values of $x$ satisfying $\displaystyle\int_0^xt^2\sin(x-t)\,dt=x^2$ in $[0,100]$.
\qn{3} Let $P(x)$ be a polynomial with real coefficients such that $\left(x^2+x+1\right)P(x-1)=\left(x^2-x+1\right)P(x)\ \forall x\in R$ and $P(1)=3$. If $\displaystyle\int_0^1\tan^{-1}\left(\dfrac{2x}{1+P\left(x^2\right)}\right)dx+\int_0^1\tan^{-1}(x+1)\,dx=\dfrac k{16}(\pi-\ln4)$, then find the value of $k$.
\qn{4} If \par\smallskip\noindent$\displaystyle\int_{-5}^{-2}\left(\dfrac{x^2-x}{x^3-3x+1}\right)^2dx+{}$\par\noindent$\displaystyle\int_{1/6}^{1/3}\left(\dfrac{x^2-x}{x^3-3x+1}\right)^2dx+{}$\par\noindent$\displaystyle\int_{6/5}^{3/2}\left(\dfrac{x^2-x}{x^3-3x+1}\right)^2dx=\dfrac pq$\par\smallskip\noindent, where $p$ and $q$ are co-primes, then $\dfrac{(6q-p)}6$ is
\qn{5} $\displaystyle\int_0^{2\pi}\max\left\{\sin x,\sin^{-1}(\sin x)\right\}dx=\dfrac{\pi^b}a-c$, then $(a+b+c)$ is $(a,b,c\in N)$
\qn{6} A college awarded 38 medals in football, 15 in basketball and 20 in cricket. If these medals went to a total of 58 men and only three men got medals in all the three sports, how many received medals in exactly two of the three sports?
\qn{7} The number of real values of $x$ satisfying the equation $3\,|x-2|+|1-5x|+4\,|3x+1|=13$ is
\qn{8} If $\displaystyle\int_0^\pi\dfrac{\cos x}{(x+4)^2}dx=P$ and $I=\displaystyle\int_0^{\pi/2}\dfrac{\sin2x}{\pi+4-2x}dx$, then the value of $P+2I-\dfrac1{\pi+4}$ is equal to $\dfrac\lambda{12}$ then $\lambda$ is
\end{multicols}
\end{document}
